\label{sec:implementation}

The \textsc{bisect-apportion} crate currently implements four divisor-method
apportionment paths in Rust:

\begin{itemize}
  \item \texttt{huntington\_hill(populations, total\_seats)}
  \item \texttt{apportionment\_divisor(populations, house\_size, RoundingRule)}
\end{itemize}

\noindent
\texttt{RoundingRule} has variants \texttt{HuntingtonHill}, \texttt{Webster},
\texttt{Adams}, and \texttt{Jefferson}.  Hamilton/largest remainder remains a
J-track mathematical and paradox-analysis reference point, but it is not yet a
public \textsc{bisect-apportion} API.

\subsection{Architecture}

All divisor methods share a priority-queue implementation:
\begin{enumerate}
  \item Assign each state 1 seat (constitutional floor).
  \item Build a max-heap of \texttt{f64} priority values.
  \item For each of the remaining $H - n$ seats, pop the state with
    highest priority, increment its seat count, and push its new
    priority back onto the heap.
\end{enumerate}
The priority formula differs by method:
for Huntington-Hill, $P_i(s) = p_i/\sqrt{s(s+1)}$.

\subsection{Current Evidence Boundary}

Earlier drafts described exact integer comparisons, a unified five-method API,
SHA-256 Census verification, and a \texttt{bisect apportion} CLI.  Those are not
the current shipped surface.  The current crate uses \texttt{f64} priority
comparisons, has 102 tests overall, and now includes a hash-bound 2020 Census
Table 1 evidence package at
\texttt{docs/examples/j-apportionment-evidence-packages/2020-census-table01/}.
That package records the Census source URL, source XLSX SHA-256, extracted
50-state apportionment rows, and verifier tests that replay Huntington-Hill
against the official 435-seat result.  Historical 2000/2010 fixtures and a
standalone \texttt{bisect apportion --verify} CLI remain future hardening work.
